\chapter{杠杆：金融世界的放大镜与绞肉机}
\chaptermeta{14}

\begin{ledetext}
同样的判断，同样的方向，同样的市场。有人赚了 40\%，有人归零。中间隔着的那个东西，你家客厅的天花板下面就有一份。
\end{ledetext}

\section{96 万，两种活法}

许工在合肥一家做汽车电子的公司当硬件工程师，太太在医院药房。2026 年，两个人凑出 96 万，包括双方父母出的一部分。

他们面前有两条路。

一条是全款买一套 96 万的小两居，从此没有月供，晚上睡得着。另一条是拿 96 万当首付，贷 224 万，买一套 320 万的房子，位置好，学区还行。

第二条路的杠杆倍数是 320 ÷ 96 = \textbf{3.33 倍}。这个数字接下来会把所有事情乘以 3.33。

\begin{tablewrap}
\begin{xltabular}{\linewidth}{>{\raggedright\arraybackslash}p{0.20\linewidth}rr}
\toprule
\thead{~} & \thead{全款买 96 万} & \thead{首付 96 万 · 贷 224 万 · 买 320 万} \\
\midrule
\textbf{自己的钱} & 96 万 & 96 万 \\
\textbf{手上握着的资产} & 96 万 & 320 万 \\
\textbf{房价涨 12\% 后的净值} & 107.5 万（+12\%） & 134.4 万（+40\%） \\
\textbf{房价跌 12\% 后的净值} & 84.5 万（-12\%） & 57.6 万（-40\%） \\
\textbf{房价跌 30\% 后的净值} & 67.2 万（-30\%） & 0 \\
\bottomrule
\end{xltabular}
\end{tablewrap}

看最后一行。

房子跌到 224 万，正好等于贷款余额。卖掉，还清银行，一分不剩。中介费和税还得自己再掏。

整个过程里，许工的判断没有错过一次，市场也没有针对他。只是他把一个 30\% 的价格波动，通过 3.33 倍的杠杆，翻译成了 100\% 的净值消失。

\begin{egbox}{还没算利息}
那 224 万不是白借的。按 2026 年 8 月 5 年期以上 LPR 3.5\% 算，30 年等额本息：

月供 \textbf{10059 元}。三十年总共还回去 362 万，其中利息 \hlmark{138 万}。

换句话说，这套 320 万的房子，许工实际要付出 96 万 + 362 万 = 458 万。房价三十年里不涨 43\%，他就是亏的。

这不是说买房不该贷款。这是说，杠杆的成本不在首付那一刻，在后面那 360 个月里，每个月 10059 块，风雨无阻。

\end{egbox}

\begin{pullquote}{}
你的房贷就是杠杆。而且对绝大多数中国家庭来说，它是这辈子唯一一笔、也是最大的一笔。
\end{pullquote}

\section{银行的 ROA 只有 0.8\%，凭什么给股东 10\%}

杠杆的正式定义很朴素：用别人的钱，控制比自己的钱更多的资产。

它对股东回报的影响可以拆成一个乘法，这个拆法一百年前是杜邦公司的会计发明的：

\begin{chainflow}
\chainnode{ROE 股东回报}\chainarrow \chainnode{ROA 资产回报}\chainarrow \chainnode{权益乘数}
\end{chainflow}

\term{权益乘数}{equity multiplier}\index{043@权益乘数}就是总资产除以自有资本。你手上有 1 块钱，撑起 12 块钱的资产，乘数就是 12。

看一家中等规模的银行长什么样：

\begin{bsheet}
  \bsside{资产 ASSETS}{%
    \bsrow{各类贷款}{682 亿}
    \bsrow{债券投资}{214 亿}
    \bsrow{现金与存放央行}{104 亿}
    \bsrowtotal{合计}{1000 亿}
  }
  \bsside{负债 + 净值}{%
    \bsrow{存款}{771 亿}
    \bsrow{同业拆借与发债}{147 亿}
    \bsrow{股东的钱}{82 亿}
    \bsrowtotal{合计}{1000 亿}
  }
\end{bsheet}

权益乘数 1000 ÷ 82 = 12.2 倍。

这家银行一年在 1000 亿资产上赚 8 亿，资产回报率 0.8\%，放在任何一个制造业老板眼里都是笑话。但对股东来说，8 亿 ÷ 82 亿 = \textbf{9.8\%}，是个体面的数字。

同一个乘数，反过来读一遍：\hlwarm{资产端只要出 82 亿的窟窿，也就是资产总额的 8.2\%，股东的钱就一分不剩。}

银行这门生意的结构就是这样。它不是因为管理不善才脆弱，它生下来就脆弱。全世界的银行监管、存款保险、最后贷款人，整套东西都是围着这 8.2\% 建的。

\begin{databox}{2026 数据}
中国商业银行净息差 2026 年二季度末为 \textbf{1.41\%}，处在历史低位。净息差被压薄，直接压的是 ROA 那一项。

ROA 掉下去，ROE 要维持住，数学上只有两条路：把权益乘数加上去，或者把成本降下去。而权益乘数是被巴塞尔协议卡死的。

我的判断是，2026 年中国银行业最真实的压力不在坏账，在息差。坏账是周期问题，息差是结构问题。

\end{databox}

\section{杀死你的不是下跌，是那通电话}

杠杆本身不杀人。

许工那套房如果跌了 30\% 他不卖，只要月供还得上，银行不会来敲门。中国的房贷极少设置盯市追加条款。他可以趴在那里等十年。

真正的绞肉机是另一套机制：\term{保证金追缴}{margin call}\index{002@保证金追缴}。

\begin{flowlist}
  \flowitem{01}{价格跌了一点}{你的仓位市值缩水，但欠的钱一分没少。}
  \flowitem{02}{担保比例跌破线}{券商或交易对手每天盯市，比例低于合同约定的数就触发条款。这是自动的，不需要任何人做判断。}
  \flowitem{03}{电话来了}{要么当天补钱，要么减仓。留给你的时间通常是一个交易日，有时候是几小时。}
  \flowitem{04}{你补不上}{能补的人早补了。被追保的时候，你身上其他资产大概率也在跌。}
  \flowitem{05}{强制平仓}{系统按市价卖，不看基本面，不等反弹，不听解释。}
  \flowitem{06}{卖压把价格再压低}{而这个更低的价格，是所有人共用的那个价格。}
  \flowitem{07}{回到 01}{另一批人的担保比例跌破线。}
\end{flowlist}

这个圈叫\term{流动性螺旋}{liquidity spiral}\index{034@流动性螺旋}。它一旦转起来，跟谁对谁错没有关系。

\begin{pullquote}{}
在这个循环里，你即使判断完全正确也会死。因为你活不到被证明正确的那一天。
\end{pullquote}

\begin{mythbox}{常见误解}
\mvfrow{误解}{"加杠杆没问题，我提前想好止损位就行了。"}
\mvfrow{事实}{止损是一份需要别人配合才能兑现的合同。你挂一个卖单，得有人接。\\ 2015 年 7 月上旬，A 股连着几天上千只股票跌停，很多人的止损单老老实实挂在那里，一天，两天，三天，成交不了。跌停板上是没有对手方的。\\ 2020 年 3 月，美股十天内四次熔断，很多流动性极好的 ETF 出现了远低于净值的成交。你以为你有一个出口，其实那扇门在最需要它的时候是从外面锁着的。\\ \textbf{止损管理的是你的纪律，管不了市场的流动性。}把仓位控制在"就算完全卖不掉也死不了"的水平，比设止损位有用得多。}
\end{mythbox}

\section{你没觉得自己借了钱的那些杠杆}

杠杆很少写着"杠杆"两个字。

\textbf{\term{期限错配}{maturity mismatch}\index{041@期限错配}。}用 7 天的钱去买 5 年的债，赚中间那段期限利差。只要每 7 天都能续上，这生意能一直做下去。整套模式的全部重量，压在"每 7 天都能续上"这一件事上。2013 年 6 月中国银行间市场的钱荒、2019 年 9 月美国回购市场那一夜隔夜利率冲到 10\%，都演过续不上是什么样子。

\textbf{\term{回购}{repo}\index{023@回购}。}拿债券做抵押借隔夜钱，借来的钱再买债，新债再拿去抵押，再借。抵押品打个 2\% 的折扣，理论上能滚出 50 倍。实践中滚不到那么高，但把自有资金放大八到十倍，在债券市场是常规操作。

\textbf{结构化产品的分层。}一个池子切成优先级和劣后级，劣后级用自己一份钱替优先级的四份钱兜底，收益也归它。它一分钱没借，但它的风险敞口是自有资金的五倍。

\textbf{期权。}一张合约对应 100 股或者一手标的，权利金只有标的价值的百分之几。买期权本身就是一种付了保险费的杠杆。

\textbf{经营杠杆。}这一种连金融都不算。一家影院，租金、折旧、最低人力是固定的，上座率从 45\% 掉到 35\%，收入掉两成，利润可能直接翻到负数。它没有任何借款，但它的成本结构就是杠杆。2020 年倒掉的那一批线下生意，很多死在这上面，不是死在负债表上。

监管机构花了很多年才认清一件事：这些形态各异的杠杆没法一个一个去堵，堵完一个还会长出下一个。只能从背面管。管不住你用多少种方式放大敞口，就管你自己垫了多少钱。

2008 年之后的巴塞尔 III，核心就三个数：资本充足率、杠杆率、流动性覆盖率。

第一个说的是你自己得垫多少钱。第二个是个不看风险权重的兜底约束，防止你把资产都塞进"看起来风险很低"的篮子里。第三个说的是：真出事的头 30 天，你手上得有足够立刻能变现的东西。

整套逻辑一句话说得完：

\begin{pullquote}{}
监管不是不让你冒险。监管是要求你冒险之前，自己先垫钱。
\end{pullquote}

为什么是资本不是存款？因为存款是别人的钱，你亏了要还；资本是股东的钱，亏了就亏了。资本这一层的厚度，决定了在损失砸到储户和纳税人身上之前，有多少缓冲。

2008 年之前，一些欧美投行的权益乘数做到 30 倍以上。3.3\% 的资产损失就能抹平全部资本。今天的大银行普遍在 10 到 15 倍。这个变化，是过去十七年金融监管做成的最实在的一件事。

\section{三倍做多，为什么长期一定亏}

这一节几乎没有科普讲清楚过，但它是纯算术，五分钟能懂。

杠杆 ETF 承诺的是\textbf{每天}跟踪标的涨跌的三倍，不是一个月三倍，不是一年三倍。为了做到每天三倍，它必须每天收盘后调仓：涨了就加仓，跌了就减仓。

这个"每天重置"的动作，在震荡市里会持续吃掉你的钱。

\begin{egbox}{两天，标的回到原点，你亏 5.45\%}

\begin{tablewrapin}
\begin{tabularx}{\linewidth}{>{\raggedright\arraybackslash}p{0.20\linewidth}rr}
\toprule
\thead{~} & \thead{标的指数} & \thead{三倍做多产品} \\
\midrule
\textbf{起点} & 100.00 & 1.0000 \\
\textbf{第一天 +10\%} & 110.00 & 1.3000 \\
\textbf{第二天 -9.09\%} & 100.00 & 0.9455 \\
\textbf{两天合计} & 一分没动 & -5.45\% \\
\bottomrule
\end{tabularx}
\end{tablewrapin}

第二天的 -9.09\% 是把 110 打回 100 需要的跌幅。三倍就是 -27.27\%，1.3 × 0.7273 = \textbf{0.9455}。

标的回到了出发点，你少了 5.45\%。

再来一遍，0.9455 × 0.9455 = 0.894。来十个这样的来回：\hlmark{0.571}。

标的指数在原地上下晃了十趟，一分钱没涨没跌。你的三倍产品少了 \textbf{42.9\%}。

\end{egbox}

这个损耗有个名字，叫\term{波动损耗}{volatility decay}\index{007@波动损耗}。它跟管理费无关，跟基金经理水平无关，它是每日重置这个机制的数学后果，躲不掉。

标的走得越颠簸，损耗越大。单边行情里三倍产品能跑出超过三倍的收益，震荡行情里它稳定失血。

所以杠杆 ETF 的说明书上那句"适合有短期交易需求的投资者"，不是免责套话，是产品的物理属性。拿三倍产品做长期持有，你赌的不是方向，是波动率。这不构成投资建议，这是算术。

\section{三个坟头}

\textbf{1998 年，LTCM。}合伙人名单里有默顿和斯科尔斯，两个人 1997 年刚拿了诺贝尔经济学奖。他们做的是相对价值套利，赚的是两个几乎一样的债券之间那点微小价差，每一笔的利润薄得像纸，所以必须上极高的杠杆才有意义。表内杠杆二十几倍，算上衍生品的名义敞口更高。

模型说，出现 1998 年 8 月那种规模的价格偏离，概率小到几百上千年才有一次。

基金 1994 年 2 月开张。1998 年 8 月，俄罗斯违约。

四年。

\textbf{2015 年，A 股场外配资。}1 比 5 是常见的，1 比 10 也有。上证从 5178 点掉下来的时候，配资盘的平仓线写在系统里，到点就砍，不看基本面，不接受申辩。卖出压低股价，压低股价触发下一批平仓线。这就是上一节那个循环，只是跑得特别快。

\textbf{2021 年，Archegos。}Bill Hwang 用总收益互换在六七家投行分别开仓，每家投行只看得见自己那一部分。等到 ViacomCBS 增发砸了股价，第一家开始平仓，剩下几家才发现原来还有别人。两天，大约 100 亿美元的自有资金归零；瑞信在这一单上亏了 55 亿美元，是它后来一系列麻烦的开头之一。

这三件事隔了二十三年，行业不同，工具不同，剧本一模一样：\textbf{高杠杆，集中持仓，对手方看不见全貌。}

\section{2026 年的杠杆在哪儿}

今天最被盯着的那块是私募信贷。它的杠杆到底高不高，答案有点反直觉。

\begin{tablewrap}
\begin{xltabular}{\linewidth}{Xr}
\toprule
\thead{时期与主体} & \thead{典型杠杆倍数} \\
\midrule
\textbf{2007 年前后的 CDO 结构} & 10–20 倍 \\
\textbf{2007 年前后的欧美投行} & 20–30 倍 \\
\textbf{2025–2026 年的私募信贷基金} & 0–2 倍 \\
\bottomrule
\end{xltabular}
\end{tablewrap}

低得多。它的负债端也更稳定，锁定期长，不像回购那样每天都要续；层层嵌套的结构也少。这几点是中金在把它叫"灰犀牛"的同时明确指出的，跟次贷不是一回事。

真正需要看的数字在另一边。

\begin{statrow}{3}
  \statitem{5000–6000}{亿美元}{美国银行体系流向私募信贷的资金}
  \statitem{3}{\%}{占银行总资产}
  \statitem{30}{\%}{占银行资本}
\end{statrow}

3\% 听起来像个零头。30\% 就不是了。

差别在分母。资产是银行手里的东西，资本是垫在最下面、亏光就完蛋的那一层。前面那张资产负债表已经算过：资产端 8.2\% 的损失就能吃掉全部资本。

套到这里：如果这 5500 亿美元的敞口最终损失三分之一，那是 1833 亿美元，等于把整个美国银行体系的资本削掉 \textbf{10\%}。不是削掉 1\%，是 10\%。

这是杠杆这一章最该记住的读数方法。\hlwarm{看一笔风险有多大，永远除以资本，不要除以资产。}除以资产是给自己看的，除以资本才是给风险看的。

\begin{mythbox}{常见误解}
\mvfrow{误解}{"杠杆是坏东西，我这个人比较保守，我从不用杠杆。"}
\mvfrow{事实}{许工的杠杆是 3.33 倍，比大部分对冲基金的股票策略还高。\\ 房贷之所以是一种"安全"的杠杆，靠的不是倍数低，而是三件事：不盯市（跌了没人打电话）、期限极长（30 年，你有时间等）、现金流来自工资而不是资产本身（房子不涨也能还上）。这三条同时成立，3.33 倍才睡得着。\\ 只要抽掉其中任何一条，同样的倍数就会变成另一种东西。比如你还款靠的是这套房的租金，而租金正好跟房价一起跌。}
\end{mythbox}

\begin{deepbox}{进阶：可以跳过}
企业和银行也有杠杆，但读法不同。企业看净负债除以 EBITDA，4 倍以上算高；银行看权益乘数和资本充足率；基金看总敞口除以净资产。

这些指标都有一个共同的盲点：它们是快照。真实的风险是"在最坏的那一天，你的资产能卖出什么价，以及有多少债在那一天到期"。这两件事不在任何一张报表上。

\end{deepbox}

\begin{takeaway}{本章一句话}
杠杆不改变你判断的对错，它只改变你能撑多久。放大收益的那个倍数，和让你归零的那个跌幅，是同一个数的正反面。

\begin{itemize}
  \item 3.33 倍杠杆下，价格跌 30\% 你的净值就是 0。倒数关系，算得清清楚楚。
  \item ROE = ROA × 权益乘数。银行 0.8\% 的资产回报能变成 9.8\% 的股东回报，代价是 8.2\% 的资产损失就归零。
  \item 致命的从来不是杠杆，是被迫平仓。止损单在流动性枯竭时执行不了。
  \item 期限错配、回购、结构化分层、期权、经营杠杆，都是杠杆，都不叫杠杆。
  \item 三倍做多产品在震荡市里稳定失血，十个来回亏 42.9\%，标的一分没动。
  \item 衡量任何一笔风险，除以资本，不要除以资产。
\end{itemize}

\end{takeaway}

\begin{bridgebox}
\textbf{下一章}杠杆能把风险放大，另一类工具能把风险切下来卖给别人。它被骂了十七年，骂得最凶的人大多没搞清它是什么，而它最早的用途，是替一个伊利诺伊的农民保住秋天的麦子价。
\end{bridgebox}

